\section{Projenin İçerdiği Yenilik Unsuru}
\label{UcuncuBolum}

Literatürdeki çalışmaların önemli bir kısmı yabancı işaret dilleri (ASL, BSL vb.)
üzerinedir ve genellikle sunucu tarafında çalışan büyük modeller kullanır. Bu
proje doğrudan TİD üzerine kurulmuştur ve AUTSL'nin resmî bölmelerini kullanarak
tekrarlanabilir bir değerlendirme düzeni oluşturur~\cite{AUTSL2020}. Kullanılan
bölmelerde işaretçi kesişmesi bulunmadığı için elde edilen doğruluk, modelin daha
önce görmediği kişiler üzerinden ölçülmüş olur.

Yenilik iddiası modelin kendisinden çok, hattın uçtan uca ve telefonda
çalışmasındadır. Şu üç nokta öne çıkmaktadır:

\begin{itemize}
  \item Girdi olarak video karesi değil iskelet koordinatları kullanılmıştır.
        Model 2,1 milyon parametreyle sınırlı kaldığı için 10 MB'lık boyut hedefi
        rahatça sağlanmış, ayrıca ham görüntü işlenmediği için kişisel veri yükü
        azalmıştır.
  \item Nokta çıkarımı hem eğitim hem uygulama tarafında aynı MediaPipe Tasks API
        ile yapılmıştır. Bu, iki taraf arasındaki nokta sıralaması ve
        normalizasyon farkından doğabilecek hataları büyük ölçüde ortadan
        kaldırmıştır. Dönüşüm sonrası Python ve LiteRT çıktılarının yüzde 99,5
        oranında aynı olması bunu desteklemektedir.
  \item Sistem tamamen çevrimdışı çalışmaktadır; model dosyası uygulamanın içinde
        taşınmakta ve çıkarım cihazda yapılmaktadır.
\end{itemize}
